% Closure of wp3 lem:KL  [drafted v2 — author line-check required]
\closuredraft
\begin{proof}
Write $\rho_k(\theta)=\rho_{k,n}(\theta)$, $v_k(\vartheta)=v_{k,n}(\vartheta)$
and abbreviate the true values $\rho_k^0$, $v_k^0$. Taking expectations of the
per-transition Gaussian log-likelihood difference and using
$\E[(Y_{k,j+1}-\rho_k(\theta)Y_{k,j})^2\mid\mathcal F_{k,j}]
=v_k^0+(\rho_k^0-\rho_k(\theta))^2Y_{k,j}^2$,
\begin{equation}\label{eq:KL-exact}
 \E_{\vzero}\big[\elln(\vzero)-\elln(\vartheta)\big]
 =\underbrace{\frac{Nn}{2}\,\overline{\psi}_{N,n}(\vartheta)}_{\text{scale part}}
 +\underbrace{\sum_{k\le N}\frac{(\rho_k^0-\rho_k(\theta))^2}
 {2v_k(\vartheta)}\sum_{j<n}\E Y_{k,j}^2}_{\text{drift part}},
\end{equation}
where
$\overline{\psi}_{N,n}(\vartheta)
:=(Nn)^{-1}\sum_{k,j}\psi\big(v_k^0/v_k(\vartheta)\big)$ and
$\psi(x):=x-1-\log x\ge0$. Both parts are non-negative, so each lower bound
may be established separately with the other parameter unrestricted.

\emph{Scale separation \eqref{eq:KL-sigma}.}
By \eqref{eq:h-exp}, uniformly in $k\le N$ and
$\theta\in[\underline\theta,\overline\theta]$,
$h_{k,n}(\theta_0)/h_{k,n}(\theta)=1+O(\lambda_N\Delta_n)$, hence
\[
 \frac{v_k^0}{v_k(\vartheta)}
 =\frac{\sigma_0^2}{\sigma^2}\big(1+O(\lambda_N\Delta_n)\big)
 \qquad\text{uniformly.}
\]
If $|\sigma-\sigma_0|\ge\varepsilon$ then
$|\sigma_0^2/\sigma^2-1|\ge\delta_\varepsilon>0$ with
$\delta_\varepsilon$ determined by $\varepsilon$ and the compact
$[\underline\sigma,\overline\sigma]$; consequently, for $N,n$ large enough
that the $O(\lambda_N\Delta_n)$ correction is at most $\delta_\varepsilon/2$
in relative size, $|v_k^0/v_k-1|\ge\delta_\varepsilon/2$ for every $(k,j)$.
Since $\psi(x)\ge c\min\{(x-1)^2,|x-1|\}$ and $v_k^0/v_k$ ranges in a fixed
compact of $(0,\infty)$, $\psi(v_k^0/v_k)\ge c_\varepsilon$, and the scale
part of \eqref{eq:KL-exact} is at least $c_\varepsilon Nn/2$. This proves
\eqref{eq:KL-sigma} with the infimum over all
$\theta\in[\underline\theta,\overline\theta]$, because the drift part is
non-negative.

\emph{Drift separation \eqref{eq:KL-theta}.}
By the mean value theorem,
$\rho_k^0-\rho_k(\theta)
=\lambda_k\Delta_ne^{-\tilde\theta\lambda_k\Delta_n}(\theta-\theta_0)$ for
some $\tilde\theta$ between $\theta$ and $\theta_0$; since
$\lambda_k\Delta_n\le\lambda_N\Delta_n\le x_0$,
$e^{-\tilde\theta\lambda_k\Delta_n}\ge c>0$, so
$|\theta-\theta_0|\ge\varepsilon$ implies
$(\rho_k^0-\rho_k(\theta))^2\ge c\,\varepsilon^2\lambda_k^2\Delta_n^2$.
For the second moments, with $s_k^2=\sigma_0^2q_k^2/(2\theta_0\lambda_k)$,
\[
 \E Y_{k,j}^2\ge s_k^2\big(1-e^{-2\theta_0\lambda_kj\Delta_n}\big)
 \ge\tfrac12 s_k^2
 \qquad\text{whenever } j\ge j_k:=\Big\lceil\frac{\log2}{2\theta_0
 \lambda_k\Delta_n}\Big\rceil .
\]
Fix $\lambda^\ast$ with $\lambda^\ast T\ge2\log2/\theta_0$; compact
resolvent implies only finitely many eigenvalues lie below $\lambda^\ast$,
so $\sum_{k:\lambda_k\ge\lambda^\ast,\,k\le N}\lambda_k\ge\LamN-C^\ast
\ge\tfrac12\LamN$ for $N$ large. For such $k$, $j_k\le n/2$, hence
$\sum_{j<n}\E Y_{k,j}^2\ge\tfrac14 ns_k^2$. Using
$v_k(\vartheta)\le\overline\sigma^2q_k^2\Delta_n(1+o(1))
\le C\overline\sigma^2q_k^2\Delta_n$,
\[
 \text{drift part}
 \ \ge\ \sum_{\substack{k\le N\\ \lambda_k\ge\lambda^\ast}}
 \frac{c\varepsilon^2\lambda_k^2\Delta_n^2}
 {C\overline\sigma^2q_k^2\Delta_n}\cdot
 \frac{n\sigma_0^2q_k^2}{8\theta_0\lambda_k}
 \ =\ \frac{c\,\sigma_0^2}{8C\overline\sigma^2\theta_0}\,
 \varepsilon^2\,T
 \sum_{\substack{k\le N\\ \lambda_k\ge\lambda^\ast}}\lambda_k
 \ \ge\ c_\varepsilon\,\LamN
\]
for $N,n\ge N_\varepsilon$, uniformly over
$\sigma\in[\underline\sigma,\overline\sigma]$, since the scale part is
non-negative. This proves \eqref{eq:KL-theta}.
\end{proof}
